\section{Introduction}

\label{sec:intro}
%----------------------------------------------------------------------

\subsection{The Reproducibility Crisis}

In 2016, a Nature survey of 1,576 researchers found that more than 70\% had failed to reproduce another scientist's experiments, and more than half had failed to reproduce their own \cite{baker2016reproducibility}. This is not a fringe concern---it is a structural failure of the scientific enterprise.

The causes are well-understood:

\begin{enumerate}
    \item \textbf{Publication Bias.} Journals publish positive results; negative results and replications languish unpublished. The literature is a biased sample of all experiments conducted.
    \item \textbf{Methodological Opacity.} Papers describe methods in prose; the actual code, data, and parameters are often unavailable.
    \item \textbf{Perverse Incentives.} Tenure, grants, and prestige are awarded for novel findings, not for careful replication.
    \item \textbf{Data Unavailability.} Datasets are proprietary, deleted after publication, or stored in formats that cannot be re-processed.
\end{enumerate}

\subsection{The Funding Bottleneck}

Scientific funding is concentrated, slow, and conservative:

\begin{itemize}
    \item The NIH funds $\sim$20\% of grant applications. The median PI age at first R01 grant is 43.
    \item Funding cycles are 12--18 months. By the time a grant is awarded, the research landscape has shifted.
    \item Geographic concentration: 80\% of global R\&D spending occurs in 10 countries.
    \item Interdisciplinary and high-risk research is systematically underfunded because review panels favor incremental advances in established fields.
\end{itemize}

\subsection{Zoo as Research Infrastructure}

\Zoo{} Network addresses these failures not through exhortation but through infrastructure. The thesis is simple: \textbf{if the tools enforce reproducibility, transparency, and fair attribution, scientists will use them because they make science easier, not harder.}

\Zoo{} provides:

\begin{enumerate}
    \item \textbf{On-Chain Provenance.} Every dataset, experiment, and publication has an immutable timestamp and authorship record on \IChain{}.
    \item \textbf{Verifiable Computation.} Experiments run in containerized environments; execution traces are hashed and anchored on-chain.
    \item \textbf{Open Data Commons.} \DID{}-authenticated data contributions with usage receipts and citation tracking.
    \item \textbf{Community Funding.} \ZIP{}s and \DAO{} governance replace opaque grant committees with transparent, milestone-based funding.
    \item \textbf{AI-Assisted Research.} \AChain{} inference for literature review, hypothesis generation, and automated analysis.
    \item \textbf{Formal Verification.} Lean~4 proof chains for mathematical results.
\end{enumerate}

\subsection{Paper Organization}

Section~\ref{sec:capsules} introduces Research Capsules. Section~\ref{sec:zips} describes ZIPs as research proposals. Section~\ref{sec:commons} presents the Data Commons. Section~\ref{sec:ai} covers AI-assisted research. Section~\ref{sec:conservation} demonstrates conservation science applications. Section~\ref{sec:verification} details formal verification. Section~\ref{sec:economics} analyzes platform economics. Section~\ref{sec:related} surveys related work. Section~\ref{sec:conclusion} concludes.

%----------------------------------------------------------------------
